\documentclass[10pt,a4paper]{amsart}
\usepackage[margin=19mm]{geometry}
\usepackage{amsmath,amssymb,mathtools}
\usepackage[scr=rsfs,cal=euler]{mathalfa}
\usepackage{xfrac}
\usepackage[unicode,hidelinks,bookmarks=false]{hyperref}
\newcommand{\ie}{i.e.\ }
\newcommand{\cf}{cf.\ }
\newcommand{\resp}{respectively\ }
\newcommand{\Subsection}{\subsection}
\providecommand{\nobreakdash}{\nobreak\hbox{-}}
\setlength{\emergencystretch}{2em}
\multlinegap0pt
\usepackage{lmodern}
\usepackage{mathtools}
\usepackage{booktabs,array,needspace}
\newtheorem{theorem}{Theorem}[section]
\newtheorem{lemma}[theorem]{Lemma}
\newtheorem{proposition}[theorem]{Proposition}
\newtheorem{corollary}[theorem]{Corollary}
\newtheorem{remark}[theorem]{Remark}
\newcommand{\AI}{\mathsf{AI}}
\newcommand{\CAI}{\mathsf{CAI}}
\newcommand{\A}{\mathcal A_4}
\newcommand{\T}{\mathcal T}
\newcommand{\Ncal}{\mathcal N}
\newcommand{\F}{\mathbb F}
\newcommand{\eps}{\varepsilon}
\newcommand{\V}{\mathsf V}
\newcommand{\bu}{\mathbf u}
\newcommand{\bv}{\mathbf v}
\newcommand{\bw}{\mathbf w}
\newcommand{\bs}{\mathbf s}
\newcommand{\bt}{\mathbf t}
\newcommand{\bfam}{\mathbf f}
\newcommand{\bK}{\mathbf k}
\newcommand{\bR}{\mathbf r}
\DeclareMathOperator{\Aff}{Aff}
\numberwithin{equation}{section}
\allowdisplaybreaks[1]
\hypersetup{pdftitle={The variety generated by all additively idempotent semirings of order four},pdfauthor={Mengya Yue and Xiaolei Shao},pdfsubject={A structural proof of nonfinite basability}}
\begin{document}
\title[The variety generated by all additively idempotent semirings]{The variety generated by all additively idempotent semirings of order four}
\author{Mengya Yue; Xiaolei Shao}
\date{}
\maketitle
\noindent\emph{Source and licence.} The variety generated by all additively idempotent semirings of order four. Version arXiv:2609.20568v1 (CC BY 4.0). This material is distributed under the Creative Commons Attribution 4.0 International licence, http://creativecommons.org/licenses/by/4.0/, as recorded in the arXiv OAI licence metadata for this exact version and shown on the abstract page https://arxiv.org/abs/2609.20568v1. Full rights notice: licenses/arxiv-2609.20568-CC-BY-4.0.txt.\par\smallskip
\noindent\textit{Editorial scope.} Counted unit: the original \texttt{theorem} environment \texttt{thm:criterion} at source lines 368--374 together with its complete proof at lines 651--688; the delivered document keeps source lines 136--689 so that every referenced label and every citation resolves inside it. Native editable source is the paper's own concatenated TeX; the AMS wrapper is editorial formatting only. No publisher class, upstream script or unrelated artwork is redistributed.
\section{Preliminaries}\label{sec:prelim}

Fix a countably infinite alphabet $X$. Write $X^+$ and $X^*$ for the free
semigroup and free monoid on $X$, and $X_c^+$ and $X_c^*$ for their
commutative counterparts. The empty word is denoted by $\eps$. By
distributivity and additive idempotence, an ai-semiring term is represented
by a finite nonempty subset of $X^+$. We write its elements as summands:
\[
\bu=u_1+\cdots+u_m=\{u_1,\ldots,u_m\}.
\]
The set $P_f(X^+)$ of such terms is the free ai-semiring on $X$, with union
as addition and setwise concatenation as multiplication; see
\cite[Section~1]{RLLC}. Thus terms will be denoted by bold letters and
individual words by ordinary letters. In particular, $w\in\bu$ means that
$w$ is a word of $\bu$.

For a word $w$, let $c(w)$ be its content, let $\ell(w)$ be its length,
and let $m(x,w)$ be the multiplicity of $x$ in $w$. For a term $\bu$, put
$c(\bu)=\bigcup_{w\in\bu}c(w)$. A substitution is an endomorphism of
$P_f(X^+)$; its values on variables are arbitrary nonempty terms. For
commutative ai-semirings the same notation is used with $X_c^+$ in place
of $X^+$. The natural homomorphism
\[
P_f(X^+)\longrightarrow P_f(X_c^+),\qquad \bu\longmapsto\overline{\bu},
\]
forgets the order of letters in each word. Equality of terms, denoted by
$=$, is equality of the represented sets; an identity is denoted by
$\approx$.

We write $\bv\preceq\bu$ for the identity $\bu+\bv\approx\bu$.
This notation is distinct from the literal inclusion $\bv\subseteq\bu$.
If $\bu=u_1+\cdots+u_m$ and $\bv=v_1+\cdots+v_t$, then
$\bu\approx\bv$ is equivalent, within $\AI$, to the identities
\[
u_i\preceq\bv\quad(1\leq i\leq m),\qquad
v_j\preceq\bu\quad(1\leq j\leq t).
\]
Indeed, these inequalities give both $\bu+\bv\approx\bu$ and
$\bu+\bv\approx\bv$. This replacement does not introduce variables.
It therefore suffices to consider inequalities $r\preceq\bs$ with $r$ a
word and $\bs$ a term.

\Needspace{8\baselineskip}
We shall use the following form of equational deduction. Derivability in
its statement is relative to the ai-semiring identities.

\begin{lemma}\label{lem:deduction}
An identity is a consequence of a set $\Sigma$ of identities if and only if
its sides can be joined by a finite sequence of contextual substitution
instances of members of $\Sigma$, used in either direction. For an
inequality $r\preceq\bs$, such a step has the form
\begin{equation}\label{eq:deduction}
\mathbf l+\mathbf p\varphi(\bs)\mathbf q
\quad\longleftrightarrow\quad
\mathbf l+\mathbf p\varphi(\bs)\mathbf q+
\mathbf p\varphi(r)\mathbf q.
\end{equation}
The factors $\mathbf p,\mathbf q$ may be absent, and $\mathbf l$ may be
absent. In $P_f(X_c^+)$ the two multiplicative contexts combine into a
single context $\bK$.
\end{lemma}
\begin{proof}
The description by contextual instances is the usual completeness theorem
for equational logic; in the ai-semiring setting it is recorded, for
example, in~\cite[Lemma~1.1]{RLLC}. The displayed form follows by expanding
a context with one distinguished occurrence of its argument. Starting
with that occurrence, addition contributes to $\mathbf l$, and left or
right multiplication contributes to the corresponding multiplicative
context, with distributivity at each stage. Substitution of $\varphi(\bs)$
and $\varphi(\bs)+\varphi(r)$ gives~\eqref{eq:deduction}. In the
commutative case the contexts multiply to $\bK=\mathbf p\mathbf q$.
\end{proof}

An absent factor is represented by $\eps$, and an absent additive context
by the empty set. These conventions concern contexts only: neither is an
allowed value of a variable under a substitution.

We next specify four commutative ai-semirings. Let
\[
H=\{1,a,a^2,\infty\},\qquad 1<a<a^2<\infty,
\]
with maximum as addition. Set $a^0=1$, $a^3=\infty$, and define
\[
a^ia^j=a^{\min\{i+j,3\}}\qquad(0\leq i,j\leq3).
\]
Let $D$ be the two-element distributive lattice, with join as addition and
meet as multiplication. For $p\in\{2,3\}$, let $G_p$ be the flat extension
of the cyclic group $C_p=\langle g\rangle$. Thus
\[
G_p=C_p\cup\{\infty\},\qquad
u+v=\begin{cases}u,&u=v,\\ \infty,&u\ne v,\end{cases}
\]
and multiplication extends the group multiplication with $\infty$
absorbing. The element names do not add constants to the signature.

The multiplication of $H$ is associative by truncated addition of
exponents and is order preserving on the chain. Hence it preserves joins
in each argument. For $G_p$, multiplication by a group element is a
permutation of $C_p$ fixing $\infty$, so it preserves the flat addition.
Multiplication by $\infty$ is constant. These observations verify the
ai-semiring axioms. The Cayley tables of $H$ and $G_3$ are given in
Table~\ref{tab:aux}.

\begin{table}[htbp]
\caption{The auxiliary semirings $H$ and $G_3$.}\label{tab:aux}
\centering
\renewcommand{\arraystretch}{1.06}
\setlength{\tabcolsep}{6pt}
$\begin{array}{c|cccc}
+&1&a&a^2&\infty\\\hline
1&1&a&a^2&\infty\\
a&a&a&a^2&\infty\\
a^2&a^2&a^2&a^2&\infty\\
\infty&\infty&\infty&\infty&\infty
\end{array}\qquad
\begin{array}{c|cccc}
\cdot&1&a&a^2&\infty\\\hline
1&1&a&a^2&\infty\\
a&a&a^2&\infty&\infty\\
a^2&a^2&\infty&\infty&\infty\\
\infty&\infty&\infty&\infty&\infty
\end{array}$
\par\medskip
$\begin{array}{c|cccc}
+&1&g&g^2&\infty\\\hline
1&1&\infty&\infty&\infty\\
g&\infty&g&\infty&\infty\\
g^2&\infty&\infty&g^2&\infty\\
\infty&\infty&\infty&\infty&\infty
\end{array}\qquad
\begin{array}{c|cccc}
\cdot&1&g&g^2&\infty\\\hline
1&1&g&g^2&\infty\\
g&g&g^2&1&\infty\\
g^2&g^2&1&g&\infty\\
\infty&\infty&\infty&\infty&\infty
\end{array}$
\end{table}

Put
\begin{equation}\label{eq:T}
\T=\V(H,D,G_2,G_3).
\end{equation}

\begin{lemma}\label{lem:inclusion}
Every ai-semiring of order at most four belongs to $\A$. In particular,
$\T\subseteq\A\cap\CAI$.
\end{lemma}
\begin{proof}
An ai-semiring $S$ embeds in an ai-semiring with one more element by
adjoining $\omega$ and making it absorbing for both operations. The old
operations are unchanged. An associativity or distributivity instance
involving $\omega$ has value $\omega$ on both sides; the remaining
instances hold in $S$. Addition remains commutative and idempotent.
Iteration gives an embedding into a four-element ai-semiring whenever
$|S|<4$. Apply this to the four algebras in~\eqref{eq:T}.
\end{proof}

For a commutative word $w$ and a finite set $Y$ containing $c(w)$, denote
its exponent vector by
\[
\nu(w)=(m(x,w))_{x\in Y}\in\mathbb Z_{\geq0}^{Y}.
\]
Its reduction modulo a prime $p$ is denoted by $\nu_p(w)$. An affine
combination of vectors over $\F_p$ is a linear combination whose
coefficients have sum one.

\begin{lemma}\label{lem:characterization}
Let $r$ be a nonempty commutative word and
$\bu=u_0+\cdots+u_t$ a commutative term.
\begin{enumerate}
\item The inequality $r\preceq\bu$ holds in $D$ if and only if
$c(u_i)\subseteq c(r)$ for some $i$.
\item For $p\in\{2,3\}$, the inequality $r\preceq\bu$ holds in $G_p$
if and only if $c(r)\subseteq c(\bu)$ and
\begin{equation}\label{eq:affine}
\nu_p(r)\in\Aff_{\F_p}\{\nu_p(u_0),\ldots,\nu_p(u_t)\}.
\end{equation}
\end{enumerate}
\end{lemma}
\begin{proof}
In $D$, a word has value one precisely when all its variables have value
one. The stated content inclusion is therefore sufficient. For necessity,
assign one to the variables of $r$ and zero to the others.

For $G_p$, a variable of $r$ outside $c(\bu)$ can be assigned $\infty$
while all other variables receive $1$, disproving the inequality. Thus
content inclusion is necessary. Choose a finite alphabet $Y$ containing
both contents. A valuation in the group part is specified by
$x\mapsto g^{\alpha_x}$, where $\alpha\in\F_p^Y$, and the value of $w$
is $g^{\nu_p(w)\cdot\alpha}$.

Let $M$ be the matrix with rows $\nu_p(u_i)-\nu_p(u_0)$,
$1\leq i\leq t$. If $\alpha\in\ker M$, all words of $\bu$ have the
same group value. Distinct group elements are incomparable in the additive
order, so validity forces
\[
(\nu_p(r)-\nu_p(u_0))\cdot\alpha=0\qquad(\alpha\in\ker M).
\]
The orthogonal complement of $\ker M$ is the row space of $M$, by
rank--nullity. This yields~\eqref{eq:affine}.

Conversely, assume the two conditions. Under a valuation for which
$\bu=\infty$, absorption is automatic. Otherwise all words $u_i$ take a
common group value $g^b$, and every variable of $\bu$, hence of $r$,
has a group value. Write $\nu_p(r)=\sum_i\lambda_i\nu_p(u_i)$ with
$\sum_i\lambda_i=1$. The exponent of the value of $r$ is
$\sum_i\lambda_i b=b$. Thus $r$ has the same value as $\bu$.
\end{proof}

\section{A sufficient condition for the nonfinite basis property}
\label{sec:criterion}

Let
\[
\Ncal=\{n\in\mathbb N:n\geq7,\ n\equiv1\pmod3\}.
\]
For $n\in\Ncal$, use pairwise distinct variables
$x_1,\ldots,x_n,y_1,\ldots,y_n,z$, with the subscripts on the $y$-variables
read modulo $n$. Define
\begin{align}
p_n&=x_1^2x_2^2\cdots x_n^2,& a_n&=x_1^2z,& q_n&=p_nz,
\label{eq:words}\\
\bfam_n&=p_n+a_n+\sum_{i=1}^n y_i y_{i+1}y_{i+2}x_i^2z.
\label{eq:family}
\end{align}
Here the products have the displayed order; no commutativity is assumed.
Let $\sigma_n$ denote the identity
\begin{equation}\label{eq:sigma}
q_n\preceq\bfam_n.
\end{equation}


\begin{theorem}\label{thm:criterion}
Let $\mathcal V$ be an ai-semiring variety containing $H,D,G_2$, and
$G_3$. If $\mathcal V$ satisfies $\sigma_n$ for infinitely many
$n\in\Ncal$, then $\mathcal V$ has no equational basis whose identities
involve a bounded number of variables. In particular, $\mathcal V$ is
nonfinitely based.
\end{theorem}


To prove the theorem we work with commutative terms. In this setting the
symbols $p_n,a_n,q_n$ denote the commutative images of the words
in~\eqref{eq:words}. Put
\begin{equation}\label{eq:u}
b_i=x_i^2y_i y_{i+1}y_{i+2}z,\qquad
\bu_n=p_n+a_n+\sum_{i=1}^n b_i=\overline{\bfam_n}.
\end{equation}
The words in $\bu_n+q_n$ are pairwise distinct. We shall characterize the
terms equivalent to $\bu_n$ in $\T$ and then examine substitutions into
these terms.

\begin{lemma}\label{lem:validT}
The variety $\T$ satisfies $q_n\preceq\bu_n$ for every $n\in\Ncal$.
\end{lemma}
\begin{proof}
In $D$, $q_n=p_nz\leq p_n$. For a valuation in $H$, if $z=1$ then
$q_n=p_n$, and if all $x_i=1$ then $q_n=a_n=z$. Otherwise some
$x_i\geq a$ and $z\geq a$, whence $x_i^2z=\infty$ and $b_i=\infty$.
This proves absorption in $H$.

Consider a valuation in $G_p$. Unless $\bu_n=\infty$, all its words
have a common group value $c$. Every variable occurs in $\bu_n$, so all
variables then take group values. When $p=2$, the squares in $p_n$ give
$p_n=1$, hence $c=1$; the equality $a_n=c$ gives $z=1$, and $q_n=c$.
When $p=3$, multiplying the equalities $b_i=c$ gives
\[
c^n=\prod_{i=1}^n b_i
=p_n\left(\prod_{i=1}^n y_i^3\right)z^n=p_nz^n.
\]
Since $n\equiv1\pmod3$ and $p_n=c$, it follows that $c=cz$. Cancellation
in the group yields $z=1$, and again $q_n=c$. Thus the inequality holds
in every generator of $\T$.
\end{proof}

\begin{lemma}\label{lem:word}
For $n\in\Ncal$ and $r\in X_c^+$,
\[
\T\models r\preceq\bu_n
\quad\Longleftrightarrow\quad
r\in\bu_n+q_n.
\]
\end{lemma}
\begin{proof}
The reverse implication follows from Lemma~\ref{lem:validT}. Suppose that
$\T\models r\preceq\bu_n$. Assigning $a$ to a single variable and $1$
to the others in $H$ gives
\begin{equation}\label{eq:degrees}
\begin{gathered}
c(r)\subseteq\{x_1,\ldots,x_n,y_1,\ldots,y_n,z\},\\
m(x_i,r)\leq2,\qquad m(y_j,r)\leq1,\qquad m(z,r)\leq1.
\end{gathered}
\end{equation}
Indeed, $\bu_n$ has value $a^2$ when the distinguished variable is $x_i$,
value $a$ when it is $y_j$ or $z$, and value $1$ for a variable outside
its content. Next, assign the nonidentity group element of $G_2$ to a
single $x_i$, and $1$ to all other variables. Every word of $\bu_n$
has value $1$, so $m(x_i,r)$ is even. Hence
\begin{equation}\label{eq:even}
m(x_i,r)\in\{0,2\}.
\end{equation}

We also have the following property:
\begin{equation}\label{eq:pairs}
\text{each pair of distinct variables in $c(r)$ occurs together in a word of $\bu_n$.}
\end{equation}
To see this, the pairs not occurring together are of the form $x_i,y_j$
with $y_j\notin\{y_i,y_{i+1},y_{i+2}\}$, or are pairs of $y$-variables
not belonging to a common triple of consecutive $y$-variables. In the
first case assign $x_i=a$ and $y_j=a^2$ in $H$, with all other variables
$1$. The value of $\bu_n$ is at most $a^2$, whereas $r=\infty$ by
\eqref{eq:even}. In the second case assign $a^2$ to both $y$-variables
and $1$ to the others, obtaining the same contradiction.

By Lemma~\ref{lem:characterization}(1), some word of $\bu_n$ has content
contained in $c(r)$. We distinguish three cases.

\emph{Case 1. $c(p_n)\subseteq c(r)$.} Every $y_j$ occurs together with
only three $x$-variables in $\bu_n$. Since $n\geq7$, property
\eqref{eq:pairs} excludes all $y$-variables from $r$. Equations
\eqref{eq:degrees} and~\eqref{eq:even} now give $r=p_n$ or $r=q_n$.

\emph{Case 2. $c(b_i)\subseteq c(r)$ for some $i$.} The occurrence of
$x_i$ excludes further $y$-variables. An additional $x_j$ would have to
occur together with all of $y_i,y_{i+1},y_{i+2}$. Its three $y$-neighbours
would therefore coincide with this triple. Distinct cyclic triples have
distinct starting indices for $n\geq7$, so $j=i$. Thus $c(r)=c(b_i)$,
and the degree restrictions imply $r=b_i$.

\emph{Case 3. $c(a_n)\subseteq c(r)$.} Write
\[
m(x_i,r)=2\delta_i,\quad \delta_i\in\{0,1\},\quad\delta_1=1,
\qquad m(y_j,r)=\eta_j\in\{0,1\}.
\]
We have $m(z,r)=1$, and~\eqref{eq:pairs} gives $\eta_j=0$ for
$4\leq j\leq n$. By Lemma~\ref{lem:characterization}(2), there exist
$\lambda_P,\lambda_A,\lambda_1,\ldots,\lambda_n\in\F_3$ such that
\begin{align*}
\nu_3(r)&=\lambda_P\nu_3(p_n)+\lambda_A\nu_3(a_n)
              +\sum_{i=1}^n\lambda_i\nu_3(b_i),\\
1&=\lambda_P+\lambda_A+\sum_{i=1}^n\lambda_i.
\end{align*}
All coefficient equations in this case are over $\F_3$. The
$z$-coordinate gives $1=\lambda_A+\sum_i\lambda_i$, so
$\lambda_P=0$. The $x_i$- and $y_j$-coordinates then give
\begin{equation}\label{eq:lambda}
\lambda_i=\delta_i-\lambda_A\mathbf1_{\{i=1\}},\qquad
\eta_j=\lambda_j+\lambda_{j-1}+\lambda_{j-2}.
\end{equation}
Consequently,
\[
\delta_j+\delta_{j-1}+\delta_{j-2}=0\quad\text{in }\F_3
\qquad(4\leq j\leq n).
\]
Each $\delta_i$ is the integer zero or one, so three such entries have
sum zero modulo three only when they are all equal. The overlapping
triples imply
\[
\delta_2=\delta_3=\cdots=\delta_n.
\]
If these entries are one, Case~1 gives $r=q_n$. If they are zero,
\eqref{eq:lambda} gives $\lambda_i=0$ for $i\ne1$ and
\[
\eta_1=\eta_2=\eta_3=1-\lambda_A.
\]
These three degrees are all zero or all one. Hence $r=a_n$ or $r=b_1$.
The three cases exhaust the content inclusion supplied by $D$.
\end{proof}

\begin{lemma}\label{lem:class}
For every $n\in\Ncal$ and every commutative term $\bv$,
\[
\T\models\bv\approx\bu_n
\quad\Longleftrightarrow\quad
\bv=\bu_n\ \text{or}\ \bv=\bu_n+q_n.
\]
\end{lemma}
\begin{proof}
Assume $\T\models\bv\approx\bu_n$. Every word of $\bv$ is absorbed
by $\bu_n$, so Lemma~\ref{lem:word} gives
$\bv\subseteq\bu_n+q_n$. We show that no word of $\bu_n$ can be
omitted from $\bv$.

For $p_n$, evaluate in $D$ with all $x_i=1$ and all $y_i=z=0$.
Then $p_n=1$ and every other word of $\bu_n+q_n$ has value zero.
For $a_n$, evaluate in $D$ with $x_1=z=1$ and all other variables zero.
Again the selected word has value one and all the others have value zero.
Finally, for $b_i$, evaluate in $H$ with $y_i=y_{i+2}=a$ and every other
variable $1$. Only the triple starting at $i$ contains both of these
$y$-variables, since $n\geq7$. Thus $b_i=a^2$, while $b_j\leq a$
for $j\ne i$ and $p_n=a_n=q_n=1$.

Each valuation separates the selected word from the sum of all remaining
words. Since every word of $\bu_n$ must be absorbed by $\bv$, all must
occur in $\bv$. Therefore
$\bu_n\subseteq\bv\subseteq\bu_n+q_n$, leaving exactly the two stated
terms. The converse is Lemma~\ref{lem:validT}.
\end{proof}

The next lemma is the main restriction on substitutions. The assumption
on $h\varphi(\bs)$ is literal inclusion of sets of commutative words.

\begin{lemma}\label{lem:substitution}
Let $n\in\Ncal$, let $\bs\in P_f(X_c^+)$ and $r\in X_c^+$, and
suppose that $\T\models r\preceq\bs$. If a substitution $\varphi$ and
a word $h\in X_c^*$ satisfy
\begin{equation}\label{eq:hyp}
h\varphi(\bs)\subseteq\bu_n,\qquad q_n\in h\varphi(r),
\end{equation}
then $|c(r)|\geq n+1$.
\end{lemma}
\begin{proof}
Each word of $h\varphi(r)$ is absorbed by $\bu_n$ in $\T$, because
$r\preceq\bs$ and $h\varphi(\bs)\subseteq\bu_n$. By
Lemma~\ref{lem:word},
\begin{equation}\label{eq:image}
h\varphi(r)\subseteq\bu_n+q_n.
\end{equation}
Also $c(r)\subseteq c(\bs)$: otherwise, in $H$, assigning $a$ to a
variable of $c(r)\setminus c(\bs)$ and $1$ to all others contradicts
$r\preceq\bs$. We use $v$ for the variables of $r$ and $\bs$.

\emph{Claim 1.} There are nonempty words $w_v\in\varphi(v)$ such that
\begin{equation}\label{eq:factor}
q_n=h\prod_{v\in c(r)}w_v^{m_v},\qquad
m_v=m(v,r)\in\{1,2\}.
\end{equation}

First, $m_v\leq2$. If $m_v\geq3$, choose the same nonempty word of
$\varphi(v)$ at all occurrences of $v$ in an expansion of
$h\varphi(r)$. Some letter would then have multiplicity at least three,
contrary to~\eqref{eq:image}.

Choose an expansion giving $q_n$. All selected words and $h$ involve
only $x_1,\ldots,x_n,z$. Suppose that the two occurrences of a variable
of multiplicity two select $w$ and $w'$. Let $k$ be the product of the
remaining selected factors, including $h$. Then $kww'=q_n$. Both
$kw^2$ and $k(w')^2$ also occur in $h\varphi(r)$. Neither $w$ nor $w'$
can contain $z$, for its diagonal product would contain $z^2$. Thus $k$
contains $z$. The diagonal products contain no $y$-variables, so by
\eqref{eq:image} each is $a_n$ or $q_n$. In integer exponent vectors,
\[
2\nu(q_n)=\nu(kw^2)+\nu(k(w')^2).
\]
Since $a_n$ is a proper divisor of $q_n$, both diagonal products must
be $q_n$. It follows that $2\nu(w)=2\nu(w')$, whence $w=w'$.
Applying this to each repeated variable proves the claim. All cancellations
here are in the free commutative monoid.

\emph{Claim 2.} The word $h$ does not contain $z$.

Suppose that $z\in c(h)$. Each $w_v$ in~\eqref{eq:factor} is then a
nonempty word in the $x$-variables. Since $c(r)\subseteq c(\bs)$, for
each $v\in c(r)$ we can choose an occurrence of $v$ in a word of $\bs$
and select $w_v$ at that occurrence. The resulting expansion, after
multiplication by $h$, belongs to $\bu_n$ and is divisible by $hw_v$.
Every word of $\bu_n$ containing $z$ has total $x$-degree two and
contains only one distinct $x_i$. Consequently, $h$ contains at most one
distinct $x_i$, and the same is true of each $w_v$.

For any $t\in\bs$, select $w_v$ at every occurrence of each
$v\in c(r)$, and arbitrary words at the other occurrences. Every selected
occurrence from $c(r)$ contributes at least one $x$-letter. Since the
expanded word, multiplied by $h$, is a word of $\bu_n$ containing $z$,
\begin{equation}\label{eq:sumdegree}
\sum_{v\in c(r)}m(v,t)\leq2\qquad(t\in\bs).
\end{equation}
Now evaluate the formal variables in $H$, assigning $a$ to all variables
of $r$ and $1$ to the others. Equation~\eqref{eq:sumdegree} gives
$\bs\leq a^2$. The inequality $r\preceq\bs$ implies $\ell(r)\leq2$.
The factorization~\eqref{eq:factor} then supplies at most
$1+\ell(r)\leq3$ distinct $x$-variables, contrary to the $n\geq7$
distinct $x$-variables of $q_n$. This proves Claim~2.

There is consequently a unique $v_0\in c(r)$ whose selected word
contains $z$, and
\begin{equation}\label{eq:v0}
m_{v_0}=1,\qquad m(z,w_{v_0})=1.
\end{equation}
Every other selected word is a nonempty word in the $x$-variables.
For every $t\in\bs$, the first condition of~\eqref{eq:hyp} implies
\begin{equation}\label{eq:sdegree}
m(v_0,t)\leq1,\qquad m(v,t)\leq2
\quad(v\in c(r)\setminus\{v_0\}).
\end{equation}
Indeed, two choices of $w_{v_0}$ would produce $z^2$, and three choices
of any other $w_v$ would produce a letter of multiplicity at least three.

There is at least one variable in $c(r)\setminus\{v_0\}$. Otherwise
$r=v_0$ and $hw_{v_0}=q_n$. An occurrence of $v_0$ in $\bs$ would give
a word of $\bu_n$ divisible by $q_n$, which is impossible.

\emph{Claim 3.} We have $h=\eps$, $w_{v_0}=z$, and each
$w_v$ with $v\ne v_0$ contains only one distinct $x$-variable.

Fix $v\in c(r)\setminus\{v_0\}$. Assign $v_0=a^2$, $v=a$ and all
other formal variables $1$ in $H$. Then $r=\infty$, so $\bs=\infty$.
By~\eqref{eq:sdegree}, this is possible only if some word of $\bs$
contains both $v_0$ and $v$. Select their words $w_{v_0}$ and $w_v$ in
an expansion. We obtain
\begin{equation}\label{eq:divides}
hw_{v_0}w_v\text{ divides a word of $\bu_n$ containing $z$.}
\end{equation}
Such a word contains only one distinct $x$-variable. If $hw_{v_0}$
contained $x_j$, equation~\eqref{eq:divides} would force every other
$w_v$ to involve only $x_j$. This contradicts~\eqref{eq:factor}.
Thus $hw_{v_0}$ contains no $x$-variable. Claim~2 and~\eqref{eq:v0}
now give $h=\eps$ and $w_{v_0}=z$. Equation~\eqref{eq:divides} also
shows that every other $w_v$ contains only one distinct $x$-variable,
proving Claim~3.

To supply all $x_1,\ldots,x_n$ in~\eqref{eq:factor}, at least $n$
distinct variables other than $v_0$ are therefore required. Hence
$|c(r)|\geq n+1$.
\end{proof}


\begin{proof}[Proof of Theorem~\ref{thm:criterion}]
Suppose that $\mathcal V$ has a basis $\Sigma$ with a variable bound $k$.
Choose $n\in\Ncal$ such that $n+1>k$ and $\mathcal V\models\sigma_n$.
The identity $\sigma_n$ is then a consequence of $\Sigma$.

Adjoin $xy\approx yx$ to the deduction system and pass to commutative
terms. This is legitimate for testing nonderivability: a deduction from
$\Sigma$ remains a deduction in the enlarged system. Moreover, all its
identities are valid in $\T$. Normalize each member of $\Sigma$ and
replace it by inequalities $r\preceq\bs$ as in Section~\ref{sec:prelim}.
The resulting set $\Sigma'$ has the same consequences modulo $\CAI$
and still uses at most $k$ variables in each identity.

By Lemma~\ref{lem:deduction}, there is a sequence of commutative terms
\[
\bu_n=\bt_0,\bt_1,\ldots,\bt_m=\bu_n+q_n
\]
whose steps are contextual instances of members of $\Sigma'$. Every
$\bt_i$ is equivalent to $\bu_n$ in $\T$. Lemma~\ref{lem:class} shows
that each is either $\bu_n$ or $\bu_n+q_n$.

Consider the first step that changes $\bu_n$. It must go from
$\bu_n$ to $\bu_n+q_n$. A contraction in~\eqref{eq:deduction} cannot
introduce a word, so the step has the form
\[
\mathbf l+\bK\varphi(\bs)
\ \longrightarrow\
\mathbf l+\bK\varphi(\bs)+\bK\varphi(r)
\]
for some $r\preceq\bs$ in $\Sigma'$. Since $q_n$ occurs only on the
right, choose $h\in\bK$ such that $q_n\in h\varphi(r)$. If the
multiplicative context is absent, take $h=\eps$. The left side gives
$h\varphi(\bs)\subseteq\bu_n$. Lemma~\ref{lem:substitution} yields
$|c(r)|\geq n+1>k$, a contradiction.

Thus there is no equational basis with a bounded number of variables.
Every finite basis has such a bound, proving nonfinite basability.
\end{proof}

\begin{thebibliography}{99}
\bibitem[RLLC]{RLLC} M. Ren, J. Liu, L. Zeng, and M. Chen, The finite basis problem for additively idempotent semirings of order four, I, \emph{Semigroup Forum} \textbf{110} (2025), 422--457. \href{https://doi.org/10.1007/s00233-025-10520-7}{doi:10.1007/s00233-025-10520-7}.
\end{thebibliography}
\end{document}
